\def\level{Première}  % Classe à droite
\def\course{NSI}
\def\chaptername{Algèbre de Boole} % Titre au centre
\def\docname{RD02-S}        % Identifiant à gauche

\pagestyle{cours_FP}
\makeheader{} % Affiche le bandeau


\begin{easybox}[Variable booléenne]{}
	\begin{itemize}
		\item     une variable booléenne est une variable qui ne peut prendre que deux valeurs : \textbf{vrai} ou \textbf{faux}.
		\item     les variables booléennes sont utilisées pour modéliser des situations à deux états, par exemple : allumé/éteint, ouvert/fermé, 1/0, \ldots{}
	\end{itemize}
\end{easybox}

\begin{easybox}[Opérateurs booléens]{}


\begin{minipage}[t]{6cm}
    \begin{center}Table de vérité du \verb+not+\end{center}
    \begin{longtable}[]{@{}cc@{}}
        \hline Expression & Résultat\tabularnewline\hline\endhead{}
        \texttt{not} True & False\tabularnewline{}
        \texttt{not} False & True\tabularnewline\hline{}
    \end{longtable}
\end{minipage}
\begin{minipage}[t]{6cm}
    \begin{center}Table de vérité du \verb+and+\end{center}
    \begin{longtable}[]{@{}cc@{}}
        \hline Expression & Résultat\tabularnewline\hline\endhead{}
        True \texttt{and} True & True\tabularnewline{}
        True \texttt{and} False & False\tabularnewline{}
        False \texttt{and} True & False\tabularnewline{}
        False \texttt{and} False & False\tabularnewline\hline{}
    \end{longtable}
\end{minipage}
\begin{minipage}[t]{6cm}
    \begin{center}Table de vérité du \verb+or+\end{center}
    \begin{longtable}[]{@{}cc@{}}
        \hline Expression & Résultat\tabularnewline\hline\endhead{}
        True \texttt{or} True & True\tabularnewline{}
        True \texttt{or} False & True\tabularnewline{}
        False \texttt{or} True & True\tabularnewline{}
        False \texttt{or} False & False\tabularnewline\hline{}
    \end{longtable}
\end{minipage}
\end{easybox}


\begin{easybox}[Element neutre]{}
	\begin{itemize}
		\item L'élément neutre de l'addition est 0: $A + 0 = A$
		\item L'élément neutre de la multiplication est 1: $A \cdot 1 = A$
	\end{itemize}
\end{easybox}

\begin{easybox}[Element absorbant]{}
	\begin{itemize}
		\item L'élément absorbant de l'addition est 1: $A + 1 = 1$
		\item L'élément absorbant de la multiplication est 0: $A \cdot 0 = 0$
	\end{itemize}
\end{easybox}

\begin{easybox}[Elements complémentaires et incompatibles]{}
	\begin{itemize}
		\item $A + \overline{A} = 1$
		\item $A \cdot \overline{A} = 0$
	\end{itemize}
\end{easybox}

\begin{propriete}[Loi de Morgan]{}
	\begin{itemize}
		\item $\overline{A + B} = \overline{A} \cdot \overline{B}$
		\item $\overline{A \cdot B} = \overline{A} + \overline{B}$
	\end{itemize}
\end{propriete}
